%----------------------------------------------------------------------------------------
%------------------------------Automatización-Industrial----------------------------------
%----------------------------------------------------------------------------------------
\section{Automatización industrial}

\subsection{Definición de automatización}
La automatización industrial se define como la aplicación de tecnologías, sistemas de control y dispositivos electrónicos para operar y supervisar procesos industriales con la mínima intervención humana. Este concepto implica la integración de componentes mecánicos, eléctricos y computacionales con el objetivo de mejorar la eficiencia, precisión y seguridad en los sistemas productivos~\cite{groover2020automation}.

\subsection{Importancia de la automatización}
La importancia de esta rama radica en su capacidad para optimizar los procesos productivos, reducir errores humanos e incrementar la calidad del producto, al mismo tiempo que disminuye los costos operativos. Además, permite la operación continua de los sistemas, el monitoreo en tiempo real y la toma de decisiones basada en datos, lo cual resulta fundamental en entornos industriales modernos altamente competitivos~\cite{groover2020automation}.

\subsection{Aplicaciones en la industria }
La automatización industrial se encuentra presente en múltiples sectores, tales como manufactura, líneas de ensamblaje, sistemas de empaquetado, control de procesos químicos, industria alimentaria y logística. Particularmente, en sistemas de tipo pick and place, la automatización permite realizar tareas repetitivas con alta precisión y velocidad, mejorando significativamente la productividad.

\subsection{Interfaces hombre-máquina (HMI)}
Un elemento fundamental dentro de la automatización industrial son las interfaces hombre máquina (HMI). Estas interfaces permiten la interacción entre el operador y el sistema automatizado, facilitando la supervisión, control y ajuste de variables del proceso. Asimismo, proporcionan visualización de datos en tiempo real, generación de alarmas y monitoreo del estado del sistema, lo cual permite una operación más intuitiva y segura de los equipos industriales.

%----------------------------------------------------------------------------------------
%--------------------------------ROBOTICA-INDUSTRIAL-------------------------------------
%----------------------------------------------------------------------------------------

\section{Robótica industrial}
La robótica industrial constituye una de las áreas más consolidadas de la automatización moderna, debido a su capacidad para ejecutar tareas repetitivas, precisas y programables dentro de entornos estructurados. Desde el punto de vista mecánico, un robot manipulador puede entenderse como una cadena cinemática compuesta por eslabones rígidos unidos mediante articulaciones, cuya función es posicionar y orientar un efector final dentro de un espacio de trabajo determinado. En este contexto, la estructura del manipulador, los grados de libertad y la naturaleza de sus articulaciones definen directamente sus capacidades de movimiento y su aplicabilidad en procesos industriales~\cite{lynch2017modern,siciliano2010robotics}.

En la industria los robots han sido ampliamente adoptados para operaciones como ensamblaje, soldadura, empaque, inspección y manipulación de materiales, principalmente porque permiten elevar la productividad, mejorar la repetibilidad y reducir la variabilidad del proceso. Dentro de estas aplicaciones, los manipuladores de arquitectura serial y, en particular, los robots cartesianos, resultan de especial interés cuando se requiere movimiento desacoplado en ejes ortogonales, facilidad de control geométrico y buena rigidez estructural. La importancia de la robótica industrial no radica únicamente en el reemplazo de tareas manuales, sino en la posibilidad de integrar sensado, modelado cinemático y estrategias de control para lograr sistemas más flexibles y adaptables~\cite{siciliano2010robotics}.

Desde la perspectiva del modelado, la robótica industrial se apoya en la cinemática para describir la relación entre las variables articulares y la pose del efector final. Esta formulación permite establecer, por un lado, qué configuración adopta el robot para un conjunto dado de coordenadas articulares y, por otro, qué valores articulares son necesarios para alcanzar una posición y orientación deseadas. En consecuencia, conceptos como espacio de trabajo, cinemática directa y cinemática inversa constituyen bases teóricas indispensables para el diseño, análisis y programación de manipuladores industriales~\cite{murray1994robotic,lynch2017modern}.

\subsection{Espacio de trabajo}
El espacio de trabajo de un manipulador se define como el conjunto de todas las configuraciones del efector final que pueden alcanzarse mediante alguna combinación admisible de variables articulares. En forma más precisa, si la cinemática directa del robot se expresa como un mapeo entre el espacio articular y el espacio de configuraciones del efector final, entonces el espacio de trabajo corresponde a la imagen de dicho mapeo. Este concepto es importante en la planificación de tareas, ya que cualquier trayectoria o punto objetivo que se asigne al robot debe pertenecer necesariamente a la región accesible por su estructura mecánica~\cite{murray1994robotic}.

Una distinción importante dentro de este concepto es la diferencia entre espacio de trabajo alcanzable y espacio de trabajo diestro. El primero corresponde al conjunto de posiciones que el efector final puede alcanzar con al menos una orientación posible, mientras que el segundo representa aquellas posiciones en las que el robot puede además adoptar orientaciones arbitrarias. Esta diferenciación resulta especialmente relevante en aplicaciones industriales, ya que no basta con que el robot llegue a un punto; en muchas tareas de manipulación también debe orientar adecuadamente la herramienta o pieza para ejecutar la operación de forma correcta~\cite{lynch2017modern}.

El espacio de trabajo está condicionado por la geometría del manipulador, el tipo de articulaciones empleadas y los límites mecánicos de cada junta. En consecuencia, la arquitectura del robot determina la forma geométrica de la región accesible. En manipuladores cartesianos, por ejemplo, el espacio de trabajo suele aproximarse a un volumen prismático, producto del desplazamiento lineal en ejes mutuamente ortogonales. Esta es una de las características que representa una ventaja en aplicaciones de \textit{pick and place} y de empaque, ya que facilita la relación entre las coordenadas físicas del entorno y las coordenadas del sistema. Además de delimitar la región físicamente accesible, el análisis del espacio de trabajo permite identificar restricciones operativas importantes, como zonas cercanas a límites articulares o regiones donde pueden presentarse singularidades cinemáticas. Por ello, su estudio no solo es útil en la etapa de diseño mecánico, sino también en la programación y validación de trayectorias, ya que contribuye a garantizar que los movimientos del robot sean realizables, seguros y compatibles con la tarea industrial asignada~\cite{lynch2017modern,siegwart2011autonomous}.

\subsection{Cinemática directa}
La cinemática directa estudia la relación entre las variables del robot y la pose del efector final, es decir, su posición y orientación respecto a un sistema de referencia. En otras palabras, permite determinar dónde se encuentra el extremo del manipulador cuando se conocen los desplazamientos o ángulos de cada articulación. Por ello, constituye una herramienta fundamental para describir geométricamente el comportamiento del robot y verificar configuraciones de trabajo~\cite{murray1994robotic}.

En manipuladores seriales, la cinemática directa se obtiene mediante la composición sucesiva de transformaciones entre eslabones. En robots cartesianos, este análisis suele simplificarse debido a que los movimientos se realizan en ejes ortogonales, lo que facilita la correspondencia entre coordenadas articulares y coordenadas espaciales. Esta característica resulta especialmente útil en aplicaciones industriales de pick and place, donde se requiere un posicionamiento claro y predecible del efector final~\cite{siciliano2010robotics,lynch2017modern}.

\subsection{Cinemática inversa}
La cinemática inversa resuelve el problema de determinar los valores articulares necesarios para que el efector final alcance una pose deseada. A diferencia de la cinemática directa, este problema suele ser más complejo, ya que puede presentar múltiples soluciones, una sola o ninguna, dependiendo de la geometría del robot y de la ubicación del objetivo dentro del espacio de trabajo~\cite{lynch2017modern,murray1994robotic}.

En aplicaciones industriales, la cinemática inversa es esencial porque las tareas suelen definirse en coordenadas cartesianas, mientras que los actuadores operan en el espacio articular. En configuraciones simples puede resolverse de forma analítica; sin embargo, en sistemas más complejos suelen emplearse métodos numéricos basados en el jacobiano. Su estudio también permite identificar configuraciones problemáticas, como las singularidades, que pueden afectar la precisión y viabilidad del movimiento del robot~\cite{siciliano2010robotics,corke2017robotics}.

%----------------------------------------------------------------------------------------
%--------------------------------ROBOTICA-INDUSTRIAL-------------------------------------
%----------------------------------------------------------------------------------------


\section{Sistemas de manipulación}
Los sistemas de manipulación comprenden el conjunto de mecanismos, actuadores, sensores y estrategias de control destinados a sujetar, mover, orientar y colocar objetos dentro de una tarea específica. En robótica industrial, estos sistemas permiten ejecutar operaciones como transferencia de piezas, ensamblaje, clasificación y empaque, integrando el manipulador con el efector final y con la lógica del proceso automatizado~\cite{siciliano2010robotics,correll2022autonomous}.

Su importancia radica en que no solo hacen posible el movimiento físico de los objetos, sino también la repetición consistente de secuencias de trabajo dentro de entornos industriales. Por ello, el desempeño de un sistema de manipulación depende de factores como la geometría del objeto, el tipo de agarre, la trayectoria seguida y la capacidad del robot para posicionarse con precisión y repetibilidad.

\subsection{Líneas de empaquetado}
Una línea de empaquetado se define como un sistema organizado de varias operaciones en el que los productos son preparados para su presentación, protección, transporte o almacenamiento. Dentro de este entorno, las tareas de manipulación son esenciales, ya que los objetos deben ser recogidos, desplazados, ordenados y colocados de forma continua y coordinada con otros elementos del proceso, como bandas transportadoras, estaciones de inspección o módulos de clasificación~\cite{siciliano2010robotics}.

La incorporación de robots en líneas de empaquetado permite aumentar la producción, mantener uniformidad en la operación y reducir la intervención manual en tareas repetitivas. Esto resulta especialmente ventajoso cuando se requiere ejecutar movimientos constantes entre posiciones definidas, conservar un ritmo de trabajo estable y adaptar la operación a distintas configuraciones de producto dentro de un sistema automatizado~\cite{siciliano2010robotics}.

\subsection{Sistemas pick and place}
Los sistemas de \textit{pick and place} son aquellas aplicaciones de manipulación en las que un robot recoge un objeto desde un lugar y lo traslada hacia otra posición definida. Aunque se trata de una tarea básica, involucra una secuencia completa de aproximación, agarre, elevación, colocación y liberación, por lo que es una de las funciones más representativas de la robótica industrial~\cite{correll2022autonomous}.

Este tipo de sistema es ampliamente utilizado en manufactura, ensamblaje y empaque, ya que permite automatizar la transferencia de objetos con rapidez y consistencia. Su desempeño depende no solo del movimiento del manipulador, sino también de la planificación de la trayectoria y la verificación de que el objeto haya sido colocado correctamente, razón por la cual suele integrarse con sensores, lógica de control secuencial y retroalimentación~\cite{correll2022autonomous,lynch2017modern}.

\subsection{Precisión en sistemas robóticos}
La precisión en sistemas robóticos se refiere al grado de cercanía entre la posición realmente alcanzada por el robot y la posición objetivo definida para la tarea. En términos prácticos, un sistema es más preciso cuando el error entre el valor deseado y el valor ejecutado es menor, por lo que este concepto depende tanto de la estructura mecánica del robot como de sus sensores, actuadores y algoritmos de control~\cite{correll2022autonomous,siciliano2010robotics}.

La norma ISO~9283 denomina a esta propiedad \textit{exactitud de posicionamiento} y establece un procedimiento para medirla: el robot visita repetidamente una pose ordenada y la exactitud corresponde a la distancia entre esa pose y el baricentro de las posiciones alcanzadas~\cite{iso9283}. Para $n$ posiciones medidas $(x_j, y_j, z_j)$ con baricentro $(\bar{x}, \bar{y}, \bar{z})$ y pose ordenada $(x_c, y_c, z_c)$:
\begin{equation}
	AP = \sqrt{(\bar{x}-x_c)^2 + (\bar{y}-y_c)^2 + (\bar{z}-z_c)^2}.
\end{equation}

\subsection{Repetibilidad}
La repetibilidad es la capacidad de un sistema robótico para regresar consistentemente a una misma posición o ejecutar una misma trayectoria bajo condiciones iguales de operación. A diferencia de la precisión, que compara el resultado con una referencia objetivo, la repetibilidad evalúa qué tan similares son entre sí múltiples ejecuciones del mismo movimiento, por lo que un robot puede ser altamente repetible aun cuando presente cierto error absoluto respecto al punto deseado~\cite{iso9283,correll2022autonomous}.

En la norma ISO~9283, la repetibilidad de posición se expresa como el radio de una esfera centrada en el baricentro que contiene a las posiciones alcanzadas~\cite{iso9283}:
\begin{equation}
	RP = \bar{l} + 3S_l, \qquad
	l_j = \sqrt{(x_j-\bar{x})^2 + (y_j-\bar{y})^2 + (z_j-\bar{z})^2},
\end{equation}
donde $\bar{l}$ es la media de las distancias $l_j$ y $S_l$ su desviación estándar. Ambas magnitudes requieren medir la posición física del efector final con un instrumento independiente del controlador.

%----------------------------------------------------------------------------------------
%-----------------------------DISEÑO-MECANICO-DEL-SISTEMA--------------------------------
%----------------------------------------------------------------------------------------

\section{Diseño mecánico del sistema}
El diseño mecánico del sistema comprende la definición de la estructura física que permitirá al robot ejecutar sus movimientos de forma estable, precisa y repetible. En un manipulador cartesiano, esta etapa incluye la selección de la arquitectura de movimiento, los grados de libertad requeridos, los elementos estructurales que guían el desplazamiento y los componentes que soportan las cargas y reducen la fricción durante la operación. Por ello, el diseño mecánico no solo determina la forma del sistema, sino también su desempeño cinemático y su viabilidad para tareas de manipulación industrial~\cite{lynch2017modern,siciliano2010robotics}.

\subsection{Movimiento en coordenadas ortogonales}
El movimiento en coordenadas ortogonales consiste en desplazar el sistema a lo largo de ejes lineales mutuamente perpendiculares, normalmente asociados con las direcciones x, y, z. Esta disposición es característica de los robots cartesianos y simplifica tanto el análisis cinemático como la programación del movimiento, ya que existe una relación más directa entre el desplazamiento físico del efector final y las coordenadas del entorno de trabajo~\cite{siciliano2010robotics,lynch2017modern}.

Desde el punto de vista mecánico, esta arquitectura ofrece ventajas importantes en sistemas de manipulación, porque facilita trayectorias rectas, reduce la complejidad geométrica y permite delimitar un volumen de trabajo. Por ello, el movimiento ortogonal resulta especialmente adecuado en aplicaciones de empaque y transferencia de objetos, donde es importante mover el efector final entre posiciones definidas de manera clara, repetitiva y controlable~\cite{siciliano2010robotics}.

\subsection{Grados de libertad}
Los grados de libertad representan el número mínimo de coordenadas independientes necesarias para describir la configuración de un sistema mecánico o de un robot. En términos generales, indican cuántos movimientos independientes puede realizar el sistema; por ello, su definición es esencial para establecer si el manipulador cuenta con la movilidad suficiente para cumplir la tarea propuesta~\cite{lynch2017modern,correll2022autonomous}.

En un robot cartesiano destinado a tareas de pick and place, los grados de libertad suelen asociarse a desplazamientos lineales en ejes ortogonales y, en algunos casos, a movimientos adicionales de orientación o agarre. La correcta selección de estos grados de libertad permite alcanzar el espacio de trabajo requerido sin añadir complejidad innecesaria al sistema, lo cual es importante para mantener un diseño mecánico simple, robusto y funcional~\cite{lynch2017modern}.

\subsection{Ejes acerados}
Los ejes acerados son elementos cilíndricos de acero utilizados como guías lineales o superficies de deslizamiento y rodadura en mecanismos de movimiento rectilíneo. En sistemas de manipulación, su función principal es proporcionar una trayectoria rígida y geométricamente estable para el desplazamiento de carros o bujes lineales, contribuyendo a conservar la alineación del movimiento y a soportar las cargas transmitidas durante la operación. Su uso es común en máquinas y sistemas lineales porque combinan resistencia mecánica, estabilidad dimensional y buena compatibilidad con elementos rodantes~\cite{thkSTShaft,nb2018life}.

\subsection{Rodamientos lineales de bolas}
Los rodamientos lineales de bolas son elementos mecánicos diseñados para permitir el desplazamiento lineal de una carga a lo largo de un eje o guía, utilizando bolas recirculantes como elementos rodantes. Su principio de funcionamiento consiste en transformar el rozamiento de deslizamiento en rozamiento de rodadura, lo que disminuye la fricción, mejora la suavidad del movimiento y favorece trayectorias más uniformes en comparación con apoyos deslizantes convencionales~\cite{thkLinearBushing,nbCatalog176E}.

\subsection{Transmisión por faja dentada y tornillo de avance}
Los mecanismos de transmisión convierten la rotación de un motor en el desplazamiento lineal de un carro. En una faja dentada, los dientes engranan con los de una polea, por lo que la transmisión no depende de la fricción y no presenta el deslizamiento de una faja plana; el desplazamiento lineal por vuelta del motor es igual al número de dientes de la polea multiplicado por el paso de la faja. Esta solución permite recorridos largos y velocidades altas con baja inercia, aunque la elasticidad de la faja y su tensado influyen en la rigidez del eje~\cite{bolton2015mechatronics}.

En un tornillo de avance, una tuerca acoplada al carro se desplaza una distancia igual al avance de la rosca por cada vuelta del tornillo. La relación de reducción es elevada, lo que aumenta la fuerza disponible y la resolución del desplazamiento a costa de la velocidad; además, la fricción de la rosca suele impedir que la carga haga girar el tornillo, característica útil en ejes verticales que deben sostener el efector final sin energía~\cite{bolton2015mechatronics}.

%----------------------------------------------------------------------------------------
%-----------------------------ELECTRONICA-Y-SISTEMAS-EMBEBIDOS--------------------------------
%----------------------------------------------------------------------------------------

\section{Electrónica y sistemas embebidos}
La electrónica y los sistemas embebidos constituyen la base funcional del sistema robótico, ya que permiten adquirir señales del entorno, procesarlas y generar acciones sobre actuadores y periféricos en tiempo real. En un robot cartesiano automatizado, esta capa integra microcontroladores, buses de comunicación, sensores, actuadores e interfaces hombre-máquina, de modo que la estructura mecánica pueda operar de forma coordinada, programable y con capacidad de retroalimentación.

\subsection{Microcontroladores}
Los microcontroladores son dispositivos electrónicos programables que integran procesador, memoria y periféricos de entrada/salida en un solo circuito, y se utilizan para ejecutar tareas de control, adquisición de datos y comunicación dentro de sistemas embebidos. En aplicaciones robóticas, su selección depende de la capacidad de cómputo, disponibilidad de interfaces, velocidad de respuesta y compatibilidad con sensores y actuadores del sistema.

\paragraph{Arduino Portenta H7} El Arduino Portenta H7 es una plataforma orientada a las aplicaciones embebidas de alto desempeño. Su procesador principal es el STM32H747XI de doble núcleo, con un Cortex-M7 a 480 MHz y un Cortex-M4 a 240 MHz, lo que permite repartir tareas entre procesamiento de alto nivel y control en tiempo real. Esta arquitectura resulta conveniente en sistemas robóticos donde se desea separar funciones como visión, comunicaciones y control de movimiento sin recurrir a múltiples placas principales~\cite{arduino2024portenta}.

\paragraph{Portenta Machine Control} La Portenta Machine Control es una placa portadora para la Portenta H7 orientada a aplicaciones industriales. Proporciona entradas y salidas digitales de 24~V, interfaces para encoders incrementales y un transceptor serial configurable como RS232, RS422 o RS485, de modo que el microcontrolador puede conectarse directamente a controladores de motores, sensores de campo y otros equipos sin acondicionamiento adicional~\cite{arduinoMachineControl}.

\paragraph{ESP32} El ESP32 es un sistema en chip de doble núcleo Xtensa LX6 de hasta 240~MHz que integra Wi-Fi, Bluetooth clásico y de baja energía, varias UART, buses I2C y salidas PWM. Su programación habitual se realiza sobre FreeRTOS, lo que permite ejecutar tareas concurrentes, y su conectividad Bluetooth lo hace adecuado para atender mandos inalámbricos e interfaces de usuario mientras otro controlador se ocupa del movimiento~\cite{espressifEsp32}.

\subsection{Sistemas operativos de tiempo real y máquinas de estado}
Un sistema operativo de tiempo real (RTOS) permite dividir el programa de un microcontrolador en tareas independientes, cada una con su prioridad y su pila, que el planificador alterna según su disponibilidad. Esta organización evita que una operación lenta o bloqueante, como esperar la respuesta de un periférico, detenga el resto del programa. Las variables compartidas entre tareas deben protegerse mediante secciones críticas, semáforos o colas para evitar lecturas inconsistentes~\cite{barry2016freertos}.

Las máquinas de estado finito complementan esta organización: el comportamiento del sistema se describe como un conjunto de estados y de transiciones activadas por eventos o condiciones. En sistemas de automatización facilitan representar secuencias de arranque, calibración, operación y error, y asegurar que cada acción se ejecute solo cuando se cumplen sus condiciones previas.

\subsection{Comunicaciones seriales}
La comunicación serial es un método de transmisión de datos en el que la información se envía de manera secuencial, bit a bit, a través de una o más líneas físicas. En sistemas embebidos se emplea ampliamente porque reduce el número de conductores necesarios, simplifica la interconexión entre dispositivos y permite enlazar microcontroladores, sensores, memorias y periféricos dentro de arquitecturas compactas.

\paragraph{Protocolo I2C}
El bus I2C es un protocolo serial síncrono que utiliza únicamente dos líneas: SDA para datos y SCL para reloj. Según la especificación oficial de NXP, permite transferencias bidireccionales orientadas a 8 bits y soporta diferentes modos de velocidad, desde 100 kbit/s en modo estándar hasta 3.4 Mbit/s en modo de alta velocidad. Su funcionamiento se basa en una arquitectura compartida entre múltiples dispositivos sobre el mismo bus, lo que lo hace especialmente útil para integrar sensores y periféricos con pocos pines. Entre sus principales ventajas se encuentran la simplicidad del cableado, la posibilidad de conectar múltiples circuitos integrados en el mismo canal y su amplia adopción en sistemas de control y electrónica embebida~\cite{nxp2021i2c}.


\paragraph{Comunicación UART}
El transmisor-receptor asíncrono universal (\textit{Universal Asynchronous Receiver-Transmitter}, UART) es un periférico que convierte datos paralelos en una secuencia de bits para su transmisión y realiza la conversión inversa durante la recepción. Se considera una comunicación asíncrona porque no incorpora una línea de reloj compartida; por ello, los dispositivos deben acordar previamente la velocidad de transmisión, el número de bits de datos, la paridad y los bits de parada. Las señales lógicas TX (transmisión) y RX (recepción) permiten establecer un intercambio de datos entre dispositivos compatibles, siempre que también se respeten sus niveles eléctricos~\cite{pena2020uart}.

En una configuración habitual de ocho bits de datos, sin paridad y con un bit de parada, denominada 8N1, cada carácter se transmite mediante un bit de inicio, ocho bits de información y un bit de parada. Por tanto, aunque el carácter útil contenga ocho bits, se transmiten diez bits en total. Cuando el formato corresponde a 8N1, la duración ideal de transmisión de un carácter puede expresarse como:
\begin{equation}
 t_{\mathrm{caracter}}=\frac{10}{R_{\mathrm{baudios}}},
\end{equation}
donde $R_{\mathrm{baudios}}$ representa la velocidad de símbolos en baudios, equivalente a bits por segundo en esta modulación binaria. Por ejemplo, a 115200 baudios, el envío de un carácter requiere aproximadamente $86.8\,\mu\mathrm{s}$, sin incluir demoras de procesamiento. Esta relación es útil para estimar el tiempo mínimo ocupado por una trama serial, pero no equivale a la latencia total de una transacción que incluya esperas, validaciones y respuestas~\cite{pena2020uart}.

La UART organiza el envío y la recepción de caracteres, pero no establece por sí misma el inicio y final de mensajes formados por varios bytes, el significado de estos, la confirmación de comandos o los mecanismos de recuperación ante la pérdida de datos. Tales funciones requieren reglas adicionales en las capas superiores. Asimismo, UART y RS485 no deben considerarse tecnologías equivalentes: la primera define el mecanismo de serialización y temporización de los caracteres, mientras que el segundo especifica características eléctricas de la transmisión sobre un enlace diferencial~\cite{kugelstadt2021rs485,pena2020uart}.

\paragraph{Estándar de comunicación RS485}
RS485, también denominado EIA-485 o TIA-485, es un estándar de interfaz eléctrica para la transmisión de señales digitales balanceadas. Establece características de los transmisores y receptores conectados a una línea diferencial, pero no determina la estructura de los paquetes, las direcciones lógicas de los nodos ni los procedimientos de solicitud y respuesta. Debido a ello, puede emplearse como capa física de diferentes protocolos de comunicación. Una de sus principales aplicaciones se encuentra en sistemas industriales con múltiples equipos, cableado compartido y posibles fuentes de interferencia electromagnética~\cite{kugelstadt2021rs485}.

La transmisión diferencial utiliza dos conductores, habitualmente identificados mediante las señales A y B, y la información se obtiene a partir de la diferencia entre sus tensiones. El receptor no interpreta únicamente el nivel de un conductor respecto a tierra, sino la diferencia de potencial entre ambas líneas. Si una perturbación externa induce una variación semejante en los dos conductores, esta aparece principalmente como ruido de modo común y puede rechazarse dentro del rango de funcionamiento del receptor. La utilización de un par trenzado contribuye a exponer ambos conductores a condiciones electromagnéticas similares y favorece ese rechazo. No obstante, la comunicación diferencial no elimina todos los errores: una terminación inadecuada, diferencias excesivas entre potenciales de referencia o perturbaciones diferenciales pueden deteriorar el enlace~\cite{kugelstadt2021rs485}.

Desde el punto de vista de la instalación, RS485 permite topologías de bus y puede conectar múltiples transceptores. En líneas de transmisión suficientemente largas o con flancos rápidos, las discontinuidades de impedancia originan reflexiones que pueden alterar la forma de los pulsos. Para reducirlas, se utilizan resistencias de terminación en los extremos físicos del bus, usualmente próximas a $120\,\Omega$ cuando esa es la impedancia característica del cable empleado. La selección de la terminación depende de la longitud, la velocidad, el tipo de conductor y las características eléctricas de los transceptores; también debe considerarse la polarización de reposo para evitar estados indeterminados cuando ningún dispositivo transmite~\cite{kugelstadt2021rs485}.

\paragraph{Comunicación half-duplex y full-duplex}
El término \textit{half-duplex} describe una comunicación bidireccional en la que los dispositivos pueden transmitir y recibir, pero no simultáneamente a través del mismo canal compartido. En una implementación RS485 de dos conductores diferenciales, los nodos utilizan un único par para ambos sentidos y deben coordinar sus intervalos de transmisión para evitar que dos transmisores intenten controlar el bus al mismo tiempo. En cambio, una implementación \textit{full-duplex} de cuatro conductores emplea dos pares diferenciales independientes, uno para cada sentido, y permite transmisión y recepción simultáneas cuando la arquitectura de los transceptores lo admite~\cite{kugelstadt2021rs485}.

En un enlace half-duplex, la dirección del flujo suele gobernarse mediante señales de habilitación del transceptor. En dispositivos como el MAX485, la entrada DE (\textit{Driver Enable}) habilita la etapa transmisora, mientras que $\overline{\mathrm{RE}}$ (\textit{Receiver Enable}, activa en nivel bajo) controla la recepción. Estas señales pueden gobernarse conjuntamente para alternar los estados de transmisión y escucha. El programa debe contemplar tanto el tiempo requerido para transmitir el último bit como un intervalo de transición antes de permitir que el siguiente nodo ocupe el bus. De lo contrario, podrían ocurrir colisiones o pérdida de caracteres durante los cambios de turno. Por ello, la gestión del acceso al medio forma parte de las reglas del sistema y no es proporcionada automáticamente por el estándar eléctrico~\cite{kugelstadt2021rs485}.

\paragraph{Transceptores RS485 y adaptación de niveles lógicos}
Un transceptor RS485 convierte las señales digitales de transmisión y recepción de un circuito lógico en niveles diferenciales apropiados para el cable de comunicación, y viceversa. Dispositivos como el MAX485 incorporan un transmisor diferencial y un receptor diferencial, cuyas conexiones lógicas se identifican comúnmente como DI (entrada de datos) y RO (salida de datos). Esta conversión permite utilizar una UART del microcontrolador como fuente y destino de los bytes, sin conectar directamente sus niveles lógicos a los conductores del bus~\cite{kugelstadt2021rs485}.

La compatibilidad eléctrica entre el microcontrolador y el transceptor constituye una consideración independiente del estándar RS485. Cuando un circuito lógico funciona a $3.3\,\mathrm{V}$ y otro a $5\,\mathrm{V}$, no debe suponerse que todas sus salidas pueden interconectarse directamente. La adaptación de niveles debe evaluarse a partir de los valores máximos permitidos por las entradas y de los umbrales de nivel alto y bajo indicados en las hojas de datos. Un convertidor de niveles lógicos puede facilitar esta conexión, pero no debe confundirse con el transceptor diferencial, ya que ambos realizan funciones distintas.

\paragraph{Protocolos por capas y encapsulamiento de datos}
En sistemas de comunicación digital resulta conveniente distinguir la capa física, el transporte de mensajes y la interpretación de los datos de aplicación. La capa física determina cómo se representan eléctricamente los bits sobre el medio; el transporte organiza los bytes en tramas reconocibles y puede añadir información para detectar mensajes dañados o identificar intercambios; finalmente, la capa de aplicación atribuye significado a campos tales como estados, coordenadas, comandos y confirmaciones. Esta separación facilita modificar una interfaz eléctrica sin redefinir necesariamente todas las estructuras de información que utiliza el sistema~\cite{kugelstadt2021rs485,pena2020uart}.

El encapsulamiento consiste en agregar campos de control alrededor de una carga útil o \textit{payload}. Según el diseño, esos campos pueden incluir identificadores de sesión, números de secuencia, versiones, longitudes y códigos de comprobación. Al recibirse una trama, cada nivel valida la información que le corresponde antes de entregarla al siguiente. Es importante diferenciar la delimitación de tramas de la validación de integridad: saber dónde empieza y termina un mensaje no asegura que su contenido se haya recibido sin alteraciones.

\paragraph{Delimitación de tramas mediante COBS}
\textit{Consistent Overhead Byte Stuffing} (COBS) es un procedimiento de codificación reversible que transforma una secuencia arbitraria de bytes en otra que no contiene el valor hexadecimal \texttt{0x00}. Este valor puede reservarse entonces como delimitador de tramas, de modo que el receptor identifique el final de un mensaje al encontrar un byte cero. El método representa las posiciones de los ceros mediante bytes de control que indican la distancia hasta el siguiente separador implícito. Cheshire y Baker~\cite{cheshire1999cobs} explican que esta codificación presenta un sobrecosto acotado y permite establecer fronteras entre paquetes en canales que transportan flujos continuos de bytes.

COBS resulta particularmente útil cuando la carga útil puede incluir cualquier valor binario, porque un delimitador simple podría confundirse con un dato válido si este no se transforma previamente. Después de un error, el receptor puede descartar los bytes de la trama afectada y volver a sincronizarse al detectar el siguiente delimitador. Sin embargo, COBS no comprueba que los datos sean correctos: su propósito principal es la delimitación y la recuperación del alineamiento de los mensajes. Para validar su integridad debe combinarse con un procedimiento adicional, como una suma de verificación o un código de redundancia cíclica~\cite{cheshire1999cobs}.

\paragraph{Detección de errores mediante redundancia cíclica}
La comprobación de redundancia cíclica (\textit{Cyclic Redundancy Check}, CRC) es una técnica de detección de errores que calcula una secuencia de bits de comprobación a partir de los datos que se desean transmitir. Su fundamento matemático consiste en interpretar la secuencia binaria como un polinomio con coeficientes en el cuerpo binario $\mathrm{GF}(2)$ y efectuar una división módulo dos respecto de un polinomio generador. Las operaciones se realizan sin acarreo; por ello, la suma y la resta binarias son equivalentes a la operación lógica XOR. El residuo calculado se incorpora al mensaje y el receptor comprueba la consistencia de los datos recibidos mediante el mismo algoritmo~\cite{williams1993crc}.

La cantidad de bits del residuo depende de la variante seleccionada. Un CRC de ocho bits produce un campo de un byte, mientras que uno de dieciséis bits ocupa dos bytes. La capacidad para identificar patrones de error depende del polinomio y de la longitud de los mensajes protegidos, entre otros factores. Un CRC permite detectar muchas alteraciones producidas durante la transmisión o el almacenamiento, pero no garantiza la detección de absolutamente todos los errores ni constituye, por sí mismo, un mecanismo de corrección o retransmisión. Si se identifica una trama dañada, el sistema debe definir si la descarta, espera una nueva actualización o solicita su envío nuevamente~\cite{williams1993crc}.

\paragraph{CRC de ocho bits (CRC-8)}
El CRC-8 calcula un residuo de ocho bits sobre una secuencia de datos. Una configuración frecuente utiliza el polinomio generador $G(x)=x^{8}+x^{2}+x+1$, representado por \texttt{0x07} al omitir el término de mayor grado. Cuando el registro se inicializa con \texttt{0x00}, los bits se procesan sin reflexión y no se aplica una operación XOR final, el resultado coincide con el modelo catalogado como CRC-8/SMBUS. Para reproducir el mismo valor en ambos extremos, no basta con indicar que se utiliza ``CRC-8'': también deben coincidir el polinomio, el estado inicial, el orden de procesamiento de bits y el valor XOR de salida~\cite{reveng2023crc,williams1993crc}.

Debe señalarse que existen variantes con el mismo polinomio que producen resultados diferentes. Por ejemplo, el control de cabecera de ATM asociado con CRC-8/I-432-1 aplica una operación XOR final diferente. Por esta razón, la denominación CRC-8/ATM utilizada en determinados códigos o descripciones debe acompañarse de los parámetros exactos del algoritmo, a fin de evitar ambigüedades técnicas. En un protocolo de aplicación, un CRC-8 permite verificar de manera compacta campos como indicadores de estado, coordenadas o comandos, con la limitación de ofrecer un campo de comprobación de un solo byte~\cite{reveng2023crc}.

\paragraph{CRC de dieciséis bits (CRC16)}
El CRC16 amplía el campo de comprobación a dieciséis bits, aunque existen múltiples variantes con diferentes polinomios e inicializaciones. La configuración conocida como CRC-16/CCITT-FALSE utiliza el polinomio $G(x)=x^{16}+x^{12}+x^{5}+1$, representado por \texttt{0x1021}, un valor inicial \texttt{0xFFFF}, procesamiento sin reflexión y XOR final \texttt{0x0000}. Esta variante también aparece catalogada con el nombre CRC-16/IBM-3740; debe distinguirse de CRC-16/MODBUS, que utiliza parámetros diferentes~\cite{reveng2023crc}.

Cuando se agrega un CRC16 a una trama de transporte, su cálculo puede abarcar no solo la carga útil, sino también campos de control tales como el identificador de sesión y el número de solicitud. De esta manera se verifica que dichos campos lleguen sin alteraciones detectables y puedan utilizarse en la correlación de mensajes. El empleo de un CRC16 no sustituye los procedimientos de delimitación, los límites temporales ni la comprobación semántica de los datos: únicamente proporciona una comprobación de integridad basada en redundancia calculada~\cite{williams1993crc}.

\paragraph{Validación por capas y comprobación semántica}
Los mecanismos de integridad pueden aplicarse en diferentes niveles de una arquitectura de comunicación. Un CRC asociado al transporte permite verificar la trama completa, incluidos sus metadatos, mientras que otro CRC asociado a la aplicación puede proteger la estructura de datos que conserva su significado independientemente del medio de transmisión. Utilizar ambos códigos introduce comprobaciones adicionales, aunque no debe interpretarse como una garantía absoluta de ausencia de errores. Tras la validación binaria también puede realizarse una comprobación semántica para identificar valores fuera de rango, versiones incompatibles, indicadores no definidos o combinaciones de campos que carezcan de sentido para el sistema de control.

La verificación semántica no es equivalente al CRC: una trama puede ser matemáticamente íntegra y, aun así, contener un comando no permitido o coordenadas fuera del espacio de trabajo. Por este motivo, los sistemas de control robustos combinan comprobaciones de integridad con criterios de aceptación relacionados con el estado de la máquina y la validez de la información.

\paragraph{Solicitudes, respuestas y confirmaciones de recepción}
Un esquema maestro-esclavo por sondeo (\textit{polling}) organiza la comunicación de manera que un coordinador inicia cada intercambio y el otro dispositivo responde. Esta estrategia es adecuada para medios half-duplex porque define turnos de acceso al canal y disminuye el riesgo de transmisiones simultáneas. En una implementación de tipo solicitud-respuesta con una sola operación pendiente, el coordinador puede transmitir una orden, esperar la respuesta asociada y, solo después, iniciar el siguiente intercambio o declarar un vencimiento del plazo establecido.

La correlación entre solicitudes y respuestas se puede establecer mediante números de secuencia y, adicionalmente, identificadores de sesión que cambien cuando se reinicia un enlace. Tales campos permiten diferenciar una respuesta actual de otra atrasada o procedente de una sesión anterior. Asimismo, el concepto de \textit{acuse de recibo} (ACK, \textit{acknowledgement}) representa una confirmación de recepción o aceptación, según las reglas específicas del protocolo. Es necesario distinguir entre recibir bytes, validar una trama y confirmar que una orden de aplicación fue ejecutada: estos acontecimientos no son necesariamente simultáneos. También pueden definirse reglas de idempotencia para evitar repetir una acción física cuando llega nuevamente una solicitud duplicada.

\paragraph{Temporización, vencimientos y recuperación del enlace}
En una comunicación serial, la recepción de datos no equivale necesariamente a la recepción de una trama válida. Por ello, los protocolos que intervienen en el control de máquinas suelen considerar un plazo máximo de respuesta, un tiempo máximo para completar una trama fragmentada y un periodo de vigencia del enlace, calculado a partir de la última comunicación aceptada. Estas magnitudes cumplen funciones distintas: un plazo de respuesta detecta una transacción vencida; el temporizador de fragmentos evita conservar indefinidamente mensajes incompletos; y la supervisión de vigencia permite identificar una interrupción prolongada de las comunicaciones.

Cuando un sistema automatizado pierde el enlace con un dispositivo que aporta consignas o información de supervisión, puede requerirse una transición hacia un estado seguro, mediante la detención de movimientos, la cancelación de acciones pendientes y la inhibición de nuevas órdenes hasta que la comunicación se estabilice. La definición de estas medidas depende del análisis de riesgos y del funcionamiento de cada máquina. En conjunto, los mecanismos de delimitación, integridad, correlación y temporización constituyen elementos complementarios para la comunicación confiable entre sistemas embebidos.

\subsection{Sensores}
Los sensores son dispositivos que convierten una magnitud física en una señal utilizable por el sistema de control. En robótica, su función es proporcionar información sobre posición, velocidad, contacto, estado lógico o variables del entorno, permitiendo implementar retroalimentación y supervisión de la operación. La selección del sensor depende de la variable a medir, la precisión requerida y la forma en que esa señal será interpretada por el microcontrolador.

\paragraph{Encoder rotativo} El encoder rotativo es un sensor que mide rotaciones, ángulo y posición rotacional. OMRON lo describe como un dispositivo que detecta posición y velocidad al convertir desplazamientos mecánicos rotacionales en señales eléctricas; en robótica, estos sensores se usan ampliamente para retroalimentar la posición y el movimiento de ejes, motores y juntas, siendo especialmente relevantes en el control de trayectoria y velocidad~\cite{omronEncoder}.

\paragraph{Finales de carrera} Un final de carrera es un interruptor de detección diseñado para activarse cuando un elemento mecánico llega a una posición límite. Según OMRON, consiste en un interruptor básico alojado en una carcasa protectora que lo resguarda de fuerzas externas, humedad, aceite, polvo y suciedad. En sistemas robóticos, se emplea para definir referencias mecánicas, evitar sobrecarreras y aumentar la seguridad operativa de los ejes lineales o articulaciones móviles~\cite{omronLimitSwitch}.

\paragraph{Sensores analógicos y digitales} Los sensores analógicos entregan una señal continua, normalmente en forma de voltaje o corriente proporcional a la magnitud medida, mientras que los sensores digitales producen valores discretos o estados lógicos. Esta diferencia condiciona la electrónica de adquisición: los sensores analógicos requieren ADC para ser procesados por el microcontrolador, mientras que los digitales pueden conectarse directamente a entradas lógicas o a buses de comunicación, según su interfaz. En consecuencia, la elección entre ambos depende de la naturaleza de la variable y la resolución requerida~\cite{stm2018adc,nxp2021i2c}.

\subsection{Interfaces y periféricos}
Las interfaces y periféricos complementan la funcionalidad del sistema embebido al permitir interacción inalámbrica, visualización de información y configuración del equipo durante la operación. En sistemas robóticos, estos elementos mejoran tanto la supervisión local como la comunicación con el usuario o con dispositivos externos.

\paragraph{Bluetooth} Bluetooth es una tecnología de comunicación inalámbrica de corto alcance diseñada para intercambiar datos entre dispositivos electrónicos sin necesidad de conexión física por cable. Su funcionamiento se basa en la transmisión de información por radiofrecuencia en la banda de 2.4 GHz, lo que permite enlazar equipos como microcontroladores, teléfonos, computadoras, sensores y periféricos. En sistemas embebidos y robóticos, Bluetooth se utiliza principalmente para habilitar comunicación inalámbrica entre el sistema y un dispositivo externo, por ejemplo para enviar comandos, monitorear variables de operación o establecer una interfaz de control remoto~\cite{elbendary2018bluetooth}.

\paragraph{Pantallas OLED}
Las pantallas OLED emplean diodos orgánicos emisores de luz, en los cuales capas delgadas orgánicas emiten luz cuando circula corriente eléctrica. Universal Display Corporation describe los OLED como dispositivos sólidos monolíticos formados por varias capas orgánicas entre electrodos conductores. En aplicaciones embebidas, su uso es conveniente por su buen contraste, bajo espesor, bajo consumo en interfaces pequeñas y capacidad para mostrar información local del sistema, como estados, variables o menús de operación~\cite{udcOLED}.

\subsection{Actuadores}
Los actuadores son dispositivos encargados de transformar una señal de control en una acción física, como movimiento lineal o rotacional, dentro del sistema. En robótica y sistemas embebidos, su función es ejecutar las órdenes generadas por el controlador para producir desplazamiento, posicionamiento o fuerza sobre un mecanismo. Por ello, los actuadores construyen el vínculo entre la etapa electrónica y la respuesta mecánica del sistema, siendo importantes para mover ejes o posicionar elementos del manipulador.

\paragraph{Motores}
Los motores son uno de los actuadores más utilizados en automatización, ya que convierten energía eléctrica en movimiento mecánico, normalmente rotacional. En aplicaciones robóticas, este movimiento puede emplearse directamente o combinarse con mecanismos como husillos, poleas, engranajes o cremalleras para obtener desplazamientos lineales o aumentar torque y control de posición. Debido a ello, los motores son fundamentales en sistemas de manipulación, porque proporcionan la energía necesaria para mover los ejes del robot de forma controlada y repetible.

\paragraph{Motores paso a paso y controladores} Un motor paso a paso gira un ángulo fijo por cada pulso que recibe su controlador, de modo que la posición del eje puede determinarse contando los pulsos ordenados sin un sensor de posición. El controlador, o \textit{driver}, recibe una señal de pulso (STEP) y una de dirección (DIR), regula la corriente de las bobinas y puede subdividir cada paso en micropasos para suavizar el movimiento. Este esquema funciona en lazo abierto: si la carga supera el par disponible o la aceleración es excesiva, el motor pierde pasos y el controlador no lo detecta~\cite{bolton2015mechatronics,orientalSteppingMotors}.

\paragraph{Servomotores} Un servomotor de modelismo integra un motor de corriente continua, un tren de engranajes, un potenciómetro y un circuito de control en lazo cerrado. Su posición angular se ordena mediante una señal de modulación por ancho de pulso con periodo típico de 20~ms (50~Hz), en la que la duración del pulso determina el ángulo, normalmente dentro de un rango cercano a $180^\circ$. Por su tamaño y facilidad de control, se utiliza con frecuencia para accionar pinzas y orientar efectores finales ligeros~\cite{bolton2015mechatronics}.

%----------------------------------------------------------------------------------------
%-----------------------------SISTEMAS-DE-VISION-COMPUTACIONAL--------------------------------
%----------------------------------------------------------------------------------------
\section{Sistemas de visión computacional}
Los sistemas de visión computacional permiten que un robot obtenga información del entorno a partir de imágenes y la transforme en datos útiles para la toma de decisiones. Estos sistemas capturan una escena mediante una cámara, procesan la información visual y extraen características relevantes como formas, bordes, regiones, objetos o posiciones dentro del espacio de trabajo. Por ello, la visión computacional es un componente clave en sistemas robóticos que deben localizar piezas, verificar condiciones del entorno o adaptar su comportamiento según la escena observada~\cite{corke2017robotics,correll2022autonomous}.

\subsection{Definición de visión computacional}
La visión computacional puede definirse como el conjunto de métodos y algoritmos que permiten a una máquina adquirir, procesar e interpretar imágenes con el propósito de obtener información significativa del entorno. A diferencia de una simple captura visual, este campo busca convertir datos bidimensionales en descripciones útiles para medición, reconocimiento o control, razón por la cual se considera una herramienta fundamental en robótica autónoma y en automatización industrial~\cite{correll2022autonomous,corke2017robotics}.

\subsection{Procesamiento de imágenes}
El procesamiento de imágenes corresponde al conjunto de operaciones aplicadas sobre una imagen para mejorarla, transformarla o preparar su contenido para etapas posteriores de análisis. Entre estas operaciones se incluyen cambios de escala, transformaciones de color, filtrado, convolución y detección de movimiento, todas orientadas a reducir ruido, resaltar estructuras y facilitar la extracción de información útil~\cite{corke2017robotics,correll2022autonomous}.

\subsection{Detección de objetos}
La detección de objetos es el proceso mediante el cual un sistema de visión identifica la presencia de uno o más objetos dentro de una imagen y determina su ubicación en la escena. Este problema va más allá del simple reconocimiento, porque no solo requiere clasificar qué aparece en la imagen, sino también localizarlo espacialmente para que el sistema robótico pueda actuar sobre él. En aplicaciones industriales, esta capacidad es esencial para tareas de pick and place, inspección visual y clasificación automática de piezas~\cite{correll2022autonomous}.

%----------------------------------------------------------------------------------------
%-----------------------------ALGORITMOS-DE-PROCESAMIENTO-DE-IMAGEN-----------------------
%----------------------------------------------------------------------------------------

\section{Algoritmos de procesamiento de imagen}
Los algoritmos de procesamiento de imagen constituyen procedimientos computacionales destinados a transformar los datos visuales capturados por un sensor en información que pueda analizarse de manera cuantitativa. Una imagen digital puede representarse mediante una función discreta $f(x,y)$, donde las coordenadas $x$ e $y$ identifican la posición de un píxel y el valor de la función describe su intensidad o sus componentes de color. A partir de esta representación es posible realizar operaciones locales, como el filtrado de vecindades, y operaciones globales, como el análisis de histogramas, la identificación de regiones y el reconocimiento de patrones~\cite{gonzalez2018digital,szeliski2022computer}.

En visión artificial, el procesamiento no se limita a mejorar la apariencia visual de una imagen. Su objetivo es obtener información útil para una tarea específica, por ejemplo, separar un objeto del fondo, determinar su ubicación en píxeles, describir su geometría o asignarle una clase. Estas operaciones pueden organizarse como etapas sucesivas; sin embargo, los sistemas basados en aprendizaje profundo también pueden realizar varias de ellas mediante un mismo modelo de inferencia. La selección del procedimiento depende de la iluminación, las características de los objetos, las restricciones de tiempo y el tipo de resultado requerido~\cite{szeliski2022computer}.

\subsection{Técnicas de procesamiento de imagen}
Las técnicas de procesamiento de imagen comprenden operaciones de acondicionamiento y transformación que modifican los datos originales para facilitar su interpretación. Entre las más comunes se encuentran la conversión entre espacios de color, el ajuste de contraste, el suavizado, la reducción de ruido y el realce de bordes. En términos generales, una transformación puede expresarse como $g(x,y)=T\{f(x,y)\}$, donde $T$ representa una operación aplicada a la imagen de entrada y $g(x,y)$ es la imagen resultante~\cite{gonzalez2018digital}.

La conversión a escala de grises permite representar una imagen mediante un único valor de intensidad por píxel, lo que simplifica determinados análisis cuando el color no constituye la característica principal. En otros casos, resulta conveniente conservar la información cromática o transformar la imagen a espacios como HSV, que permiten separar parcialmente el tono de la intensidad luminosa. Estas representaciones pueden facilitar la distinción entre regiones cuando existe suficiente contraste respecto del fondo~\cite{szeliski2022computer}.

El filtrado espacial permite modificar el valor de un píxel a partir de una vecindad. Los filtros de promedio y gaussianos reducen variaciones locales de intensidad, aunque pueden suavizar los límites de los objetos. El filtro de mediana, por su parte, resulta útil frente a ciertos tipos de ruido impulsivo porque reemplaza el valor de un píxel por la mediana de su vecindad. La elección del filtro implica un compromiso entre la supresión de ruido y la conservación de detalles relevantes. El realce de bordes, mediante operadores basados en gradientes, busca identificar transiciones pronunciadas de intensidad asociadas con límites de regiones o estructuras de interés~\cite{gonzalez2018digital}.

Otra operación relevante es la transformación geométrica, que permite relacionar puntos de una imagen con posiciones en otro sistema de coordenadas. Entre sus aplicaciones se encuentran la rectificación de perspectiva, la corrección de distorsiones y la conversión de posiciones medidas en píxeles a coordenadas de un plano de trabajo. Cuando los puntos de interés pertenecen aproximadamente a un mismo plano, una homografía proporciona una relación proyectiva entre ambos sistemas de coordenadas, siempre que se disponga de correspondencias adecuadas para su estimación~\cite{szeliski2022computer,corke2017robotics}.

\subsection{Segmentación}
La segmentación consiste en dividir una imagen en regiones que presentan propiedades comunes o que corresponden a objetos de interés. Su finalidad es establecer qué píxeles pertenecen al fondo y cuáles a una o varias regiones relevantes. A diferencia de la detección de objetos, que puede entregar únicamente una caja delimitadora, la segmentación describe la pertenencia de los píxeles a una región. Esto permite estimar con mayor detalle el contorno, el área y la forma de los objetos~\cite{gonzalez2018digital,szeliski2022computer}.

Una técnica elemental es la umbralización, en la cual los píxeles se separan de acuerdo con su intensidad o con alguna componente de color. Para una imagen en escala de grises, una segmentación binaria puede definirse mediante un umbral $T$:
\begin{equation}
    s(x,y)=
    \begin{cases}
        1, & f(x,y)\geq T,\\
        0, & f(x,y)<T.
    \end{cases}
\end{equation}
La selección del umbral puede ser manual, adaptativa o automática. El método de Otsu determina un umbral a partir del histograma de intensidades mediante la maximización de la varianza entre clases. No obstante, la eficacia de este tipo de técnicas depende de que exista una separación suficientemente clara entre las distribuciones del objeto y del fondo~\cite{otsu1979threshold}.

También se emplean métodos basados en crecimiento de regiones, agrupamiento de características y análisis de contornos. Las operaciones morfológicas, como la erosión y la dilatación, permiten modificar regiones binarias para suprimir elementos pequeños, unir componentes próximos o mejorar la continuidad de las formas. En enfoques modernos, las redes neuronales pueden producir una etiqueta por píxel, lo que se conoce como segmentación semántica; cuando además se distinguen objetos individuales de una misma categoría, se habla de segmentación de instancias~\cite{gonzalez2018digital,szeliski2022computer}.

La segmentación resulta especialmente útil cuando se requiere conocer la forma completa de un objeto. Sin embargo, su utilización no es indispensable en todos los sistemas de reconocimiento: para determinadas tareas de localización y manipulación puede ser suficiente una caja delimitadora producida por un detector, siempre que la precisión geométrica obtenida sea adecuada para la aplicación.

\subsection{Extracción de características}
La extracción de características es el proceso de obtener descriptores numéricos o geométricos que resuman propiedades relevantes de una imagen o de una región detectada. En lugar de analizar todos los píxeles durante las etapas posteriores, se emplea una representación más compacta que puede incluir dimensiones, posición, área, perímetro, orientación, textura y relaciones de forma. La utilidad de cada descriptor depende de su capacidad para distinguir los objetos considerados y de su estabilidad frente a cambios de iluminación, escala o punto de vista~\cite{gonzalez2018digital,szeliski2022computer}.

Entre las características geométricas se encuentran el centroide y la caja delimitadora. Para una región binaria con $N$ píxeles pertenecientes al objeto, el centroide puede estimarse mediante los promedios de las coordenadas de dichos píxeles:
\begin{equation}
    \bar{x}=\frac{1}{N}\sum_{i=1}^{N}x_i,
    \qquad
    \bar{y}=\frac{1}{N}\sum_{i=1}^{N}y_i.
\end{equation}
La caja delimitadora, por otra parte, se define mediante los valores extremos de las coordenadas o mediante su centro, ancho y alto. Su relación de aspecto permite distinguir objetos predominantemente alargados en dirección horizontal o vertical, aunque no determina por sí sola la orientación exacta de una pieza irregular. Para obtener información más detallada pueden utilizarse momentos de imagen, descriptores de contorno o estimaciones de ejes principales~\cite{gonzalez2018digital}.

Las características también pueden basarse en puntos distintivos, bordes y patrones locales de textura. En los modelos de aprendizaje profundo, una parte importante de estas representaciones se obtiene de manera automática mediante capas sucesivas de procesamiento, en lugar de seleccionarse exclusivamente a mano. En aplicaciones robóticas, los descriptores geométricos permiten relacionar la salida de la visión con acciones posteriores, como el posicionamiento del efector final o la selección de una orientación de agarre. Para convertir una posición de imagen en una posición física es necesaria, además, una calibración geométrica apropiada~\cite{szeliski2022computer,corke2017robotics}.

\subsection{Clasificación de objetos}
La clasificación de objetos consiste en asignar una categoría a una imagen o a una región visual según sus características. En un enfoque convencional, se extraen descriptores y se utilizan reglas de decisión o modelos estadísticos para determinar la clase. Entre los métodos de clasificación se encuentran los árboles de decisión, las máquinas de vectores de soporte y las redes neuronales. La elección del método depende de la variabilidad de los objetos, la cantidad de datos disponibles y los recursos de procesamiento~\cite{szeliski2022computer}.

Las redes neuronales convolucionales permiten aprender representaciones jerárquicas de la información visual. Durante el entrenamiento, sus parámetros se ajustan utilizando ejemplos etiquetados para reducir una función de pérdida; durante la inferencia, el modelo procesa imágenes nuevas y estima la pertenencia a las categorías aprendidas. Su desempeño depende de factores como la representatividad del conjunto de entrenamiento, la calidad de las etiquetas y las diferencias entre las imágenes de entrenamiento y las condiciones reales de operación~\cite{gonzalez2018digital,szeliski2022computer}.

Es importante distinguir la clasificación de la detección de objetos. Un clasificador determina la categoría asociada con una entrada visual, mientras que un detector identifica una o varias instancias dentro de una imagen y proporciona información sobre su localización, generalmente mediante cajas delimitadoras y puntuaciones de confianza. Los detectores modernos pueden integrar la clasificación y la localización en una misma arquitectura. Esta combinación resulta relevante para la automatización de tareas de manipulación, en las cuales es necesario conocer tanto el tipo de pieza como su ubicación espacial~\cite{szeliski2022computer}.

La calidad de la clasificación puede evaluarse mediante matrices de confusión y métricas como precisión, exhaustividad y exactitud, dependiendo del problema. En detección también se considera el grado de coincidencia entre las cajas predichas y las de referencia. Estas medidas permiten diferenciar el desempeño del reconocimiento visual de otros factores del sistema robótico, como los errores de calibración, las demoras de comunicación y el posicionamiento mecánico.

\subsection{Detectores YOLO y segmentación de instancias}
Los detectores de una etapa de la familia YOLO (\textit{You Only Look Once}) procesan la imagen completa en una sola pasada de una red convolucional y predicen simultáneamente las cajas delimitadoras, las clases y sus confianzas, lo que permite inferencia en tiempo real~\cite{redmon2016yolo}. Las versiones recientes de la biblioteca Ultralytics, como YOLO11, ofrecen variantes para detección, cajas orientadas y segmentación de instancias; en esta última, además de la caja, la red predice una máscara con los píxeles de cada objeto. Los modelos se entrenan a partir de pesos preentrenados con un conjunto propio de imágenes etiquetadas y se exportan a formatos intermedios, como ONNX, para su ejecución en otros dispositivos~\cite{jocher2024yolo11}.

\subsection{Métricas de evaluación de modelos de detección}
La coincidencia entre una región predicha $P$ y una región de referencia $G$ se mide con la intersección sobre la unión:
\begin{equation}
	\mathrm{IoU} = \frac{\lvert P \cap G \rvert}{\lvert P \cup G \rvert}.
\end{equation}
Una predicción se considera verdadera positiva ($TP$) si su clase es correcta y su IoU supera un umbral; las predicciones restantes son falsas positivas ($FP$) y los objetos no detectados, falsos negativos ($FN$). A partir de estos conteos se definen la precisión $P = TP/(TP+FP)$ y la exhaustividad $R = TP/(TP+FN)$. La precisión promedio (AP) es el área bajo la curva precisión--exhaustividad de una clase, y el mAP es su promedio entre clases. El indicador mAP50 usa un umbral de IoU de 0,5, mientras que mAP50-95 promedia los umbrales de 0,50 a 0,95, por lo que exige una localización más ajustada~\cite{szeliski2022computer,jocher2024yolo11}. Estas métricas se calculan sobre un conjunto de validación y describen el modelo, no el desempeño completo del sistema robótico.

\subsection{Cámaras con procesamiento embebido}
Las cámaras con procesamiento embebido integran el sensor de imagen y un procesador capaz de ejecutar modelos de inferencia, de modo que entregan al microcontrolador resultados ya procesados, como clases, posiciones y dimensiones de los objetos, en lugar de imágenes completas. La HUSKYLENS 2 de DFRobot pertenece a esta categoría: ofrece algoritmos predefinidos, entre ellos el reconocimiento de marcadores AprilTag, permite cargar modelos personalizados convertidos a su formato propio y se comunica por UART o I2C~\cite{dfrobotHuskylens2,dfrobotHuskylensLib}. Esta arquitectura reduce la carga del microcontrolador, pero limita la información disponible a la que expone la interfaz de la cámara.

\subsection{Marcadores fiduciales y homografía}
Los marcadores fiduciales son patrones impresos diseñados para ser detectados e identificados de forma robusta en una imagen. Los AprilTag codifican un identificador en una matriz binaria con un borde negro y permiten obtener con precisión de subpíxel la posición de sus esquinas y de su centro, aun con cambios de iluminación y perspectiva~\cite{olson2011apriltag,wang2016apriltag2}.

Cuando los puntos observados pertenecen a un mismo plano, la relación entre las coordenadas de imagen $(u,v)$ y las coordenadas métricas del plano $(X,Y)$ es una homografía, es decir, una transformación proyectiva representada por una matriz $H$ de $3\times3$ definida salvo un factor de escala:
\begin{equation}
	\lambda
	\begin{bmatrix} X \\ Y \\ 1 \end{bmatrix}
	= H
	\begin{bmatrix} u \\ v \\ 1 \end{bmatrix}.
\end{equation}
Si se fija $h_{22}=1$, la matriz tiene ocho incógnitas y cada correspondencia entre un punto de imagen y uno del plano aporta dos ecuaciones lineales, por lo que cuatro puntos no colineales bastan para determinarla. Con más correspondencias, la matriz se estima por mínimos cuadrados. La homografía solo es válida para puntos del plano calibrado: los objetos con altura se proyectan con un error que crece con su distancia al plano~\cite{hartley2004multiple,szeliski2022computer}.

%----------------------------------------------------------------------------------------
%---------------------------------SISTEMAS-DE-CONTROL------------------------------------
%----------------------------------------------------------------------------------------

\section{Sistemas de control}
Los sistemas de control comprenden los elementos y procedimientos utilizados para influir sobre el comportamiento de un proceso con el propósito de alcanzar una condición de operación deseada. En un sistema mecatrónico, dichos elementos pueden incluir un controlador programable, sensores, actuadores y una planta física. Su diseño requiere establecer las variables que se desean regular, las señales disponibles para la medición y los efectos de las perturbaciones externas~\cite{nise2015control,astrom2020feedback}.

\subsection{Definición de sistema de control}
Un sistema de control puede describirse mediante una señal de referencia $r(t)$, que representa el comportamiento deseado, una acción de control $u(t)$, que actúa sobre la planta, y una salida $y(t)$, que describe la respuesta del proceso. El controlador determina la acción de control en función de la información disponible y de los objetivos operativos. En sistemas digitales, la adquisición de señales y el cálculo de las órdenes se realizan en instantes discretos, con un periodo de muestreo que condiciona la rapidez de respuesta y la calidad de la regulación~\cite{franklin2010feedback,nise2015control}.

Las propiedades de un sistema de control se analizan mediante criterios como estabilidad, precisión, rapidez de respuesta y rechazo de perturbaciones. La precisión se relaciona con la diferencia entre la salida obtenida y la referencia; la rapidez describe el tiempo requerido para alcanzar una condición deseada; y la estabilidad caracteriza la evolución de la respuesta ante cambios de consigna o perturbaciones. Estas propiedades no dependen solamente del algoritmo de control, sino también de la dinámica de la planta, los actuadores, los sensores y los retrasos presentes en la comunicación y el procesamiento~\cite{astrom2020feedback}.

\subsection{Sistemas de lazo abierto}
En un sistema de lazo abierto, la acción de control se establece sin utilizar una medición de la salida para corregir continuamente la diferencia respecto de la referencia. Por tanto, el controlador genera las órdenes a partir de una secuencia o un modelo previamente definidos, suponiendo que la planta responderá de manera suficientemente predecible. Esta arquitectura puede ser sencilla de implementar y adecuada cuando las perturbaciones son pequeñas y la relación entre la entrada y la salida se conoce con suficiente precisión~\cite{nise2015control}.

Un ejemplo es el accionamiento de un motor paso a paso mediante una cantidad determinada de pulsos, calculada a partir del desplazamiento solicitado y de una relación conocida de pasos por unidad de longitud. Si no existe una medición independiente que confirme la posición realmente alcanzada, el controlador supone que los pulsos ordenados se tradujeron en movimiento efectivo. En estas condiciones, fenómenos como pérdida de pasos, sobrecarga, deslizamiento o interferencias mecánicas pueden producir errores que el sistema no corrige automáticamente~\cite{franklin2010feedback,nise2015control}.

Los sistemas de lazo abierto ofrecen ventajas de simplicidad, costo y facilidad de programación, pero su desempeño depende de la repetibilidad de la planta y de las condiciones de operación. Debe distinguirse entre utilizar sensores para establecer referencias iniciales o detectar límites mecánicos y disponer de una realimentación continua de posición: la presencia de finales de carrera no convierte por sí sola a un accionamiento en un servosistema de posición en lazo cerrado.

\subsection{Sistemas de lazo cerrado}
Un sistema de lazo cerrado incorpora información de la salida medida para modificar la acción de control. En su forma más habitual, la señal de realimentación se compara con la referencia y se obtiene un error:
\begin{equation}
    e(t)=r(t)-y(t).
\end{equation}
A partir de este error, el controlador calcula una acción destinada a reducir la diferencia entre la condición deseada y la medida. La realimentación permite responder a ciertas perturbaciones y compensar, dentro de los límites del sistema, variaciones en el comportamiento de la planta~\cite{astrom2020feedback,nise2015control}.

En una implementación digital, el error puede calcularse como $e[k]=r[k]-y[k]$, donde $k$ representa el instante de muestreo. La respuesta depende de la frecuencia de adquisición, el tiempo de procesamiento, la calidad de la medición y la dinámica de los actuadores. Entre las estrategias de control más utilizadas se encuentran los controladores proporcional, proporcional-integral y proporcional-integral-derivativo; no obstante, la existencia de realimentación no implica necesariamente el empleo de un controlador PID~\cite{franklin2010feedback}.

En sistemas robóticos pueden coexistir distintos niveles de realimentación. Un encoder puede proporcionar información sobre el desplazamiento de una banda transportadora, mientras que un sensor de posición instalado directamente en un eje puede permitir corregir el movimiento de ese eje. Asimismo, un sistema de visión puede actualizar objetivos y decisiones a nivel de supervisión sin establecer necesariamente un lazo cerrado continuo de posición sobre cada motor. Por ello, al caracterizar una arquitectura de control es necesario identificar qué variable se mide y qué acción se modifica realmente a partir de esa medición.

\subsection{Ventajas del control en lazo cerrado}
La principal ventaja del control en lazo cerrado es la posibilidad de ajustar el comportamiento de la planta cuando la salida difiere de la referencia. Esto puede mejorar la precisión frente a perturbaciones, reducir el efecto de determinadas incertidumbres del modelo y aumentar la capacidad de adaptación ante variaciones de carga o de condiciones operativas. La realimentación también hace posible supervisar variables del proceso y detectar desviaciones que pasarían inadvertidas en un sistema de lazo abierto~\cite{astrom2020feedback}.

En aplicaciones de movimiento, una medición adecuada puede contribuir a corregir errores de posición o velocidad. Por ejemplo, cuando una banda transportadora cambia de velocidad, el uso de un encoder permite estimar su desplazamiento a partir de las cuentas medidas en lugar de depender exclusivamente de una velocidad supuesta constante. Sin embargo, el beneficio obtenido depende de cómo se utilice esa información: medir la banda no equivale a medir directamente la posición real del efector final ni garantiza que todos los ejes estén regulados en lazo cerrado.

El control realimentado presenta también limitaciones. Los sensores introducen ruido, incertidumbre y posibles retrasos; el procesamiento digital y las comunicaciones añaden latencia; y una selección inadecuada de ganancias o de leyes de control puede provocar oscilaciones o inestabilidad. Por tanto, la realimentación no asegura por sí sola un mejor desempeño. Sus ventajas se materializan cuando se diseñan de forma coherente el sistema de medición, el algoritmo de control, los actuadores y las condiciones de operación~\cite{franklin2010feedback,nise2015control}.
